\textbf{Phần 1. Trắc nghiệm lựa chọn}
\setcounterpageref{ex}{0}

\begin{ex} %[2H9V47]
Hình nào dưới đây có các mặt đều là hình vuông bằng nhau?
\choice
{Khối hộp chữ nhật}
{Khối trụ}
{\True Khối lập phương}
{Khối cầu}
\loigiai{
Khối lập phương có 6 mặt đều là hình vuông bằng nhau.
\textbf{Đáp án của câu hỏi là C}
}
\end{ex}

\begin{ex} %[2H9V47]
Bao diêm, viên gạch xây nhà thường có hình dạng là khối gì?
\choice
{Khối lập phương}
{\True Khối hộp chữ nhật}
{Khối trụ}
{Khối cầu}
\loigiai{
Bao diêm và viên gạch có dạng khối hộp chữ nhật.
\textbf{Đáp án của câu hỏi là B}
}
\end{ex}

\begin{ex} %[2H9V47]
Xếp $8$ khối lập phương nhỏ cạnh $1\text{ cm}$ thành một khối lập phương lớn. Khối lập phương lớn có cạnh dài bao nhiêu?
\choice
{$3\text{ cm}$}
{$4\text{ cm}$}
{$1\text{ cm}$}
{\True $2\text{ cm}$}
\loigiai{
Khối lập phương lớn gồm $2 \times 2 \times 2 = 8$ khối nhỏ, có cạnh dài $2\text{ cm}$.
\textbf{Đáp án của câu hỏi là D}
}
\end{ex}

\begin{ex} %[2H9V47]
Trong các đồ vật sau: quả bóng bàn, quả bóng chuyền, quả bóng rổ. Chúng đều có điểm chung là có dạng
\choice
{\True Khối cầu}
{Khối trụ}
{Khối tròn dẹt}
{Khối lập phương}
\loigiai{
Các quả bóng đều có dạng khối cầu.
\textbf{Đáp án của câu hỏi là A}
}
\end{ex}

\begin{ex} %[2H9V47]
Một người thợ xây dùng $45$ viên gạch hình hộp chữ nhật để xây bệ đỡ, sau đó dùng thêm $28$ viên gạch nữa. Bác thợ đã dùng tất cả bao nhiêu viên gạch?
\choice
{$71$ viên gạch}
{$72$ viên gạch}
{\True $73$ viên gạch}
{$75$ viên gạch}
\loigiai{
$45 + 28 = 73$ (viên gạch).
\textbf{Đáp án của câu hỏi là C}
}
\end{ex}

\begin{ex} %[2H9V47]
Thùng đựng nước hình trụ có thể chứa $20\text{ L}$ nước. Hiện trong thùng có $12\text{ L}$ nước. Cần đổ thêm bao nhiêu lít nước để đầy thùng?
\choice
{$7\text{ L}$}
{\True $8\text{ L}$}
{$9\text{ L}$}
{$6\text{ L}$}
\loigiai{
$20 - 12 = 8\text{ (L)}$.
\textbf{Đáp án của câu hỏi là B}
}
\end{ex}

\begin{ex} %[2H9V47]
Con xúc xắc đồ chơi để chơi cờ cá ngựa có dạng hình khối gì?
\choice
{Khối trụ}
{Khối cầu}
{Khối chữ nhật}
{\True Khối lập phương}
\loigiai{
Xúc xắc có 6 mặt vuông bằng nhau nên có dạng khối lập phương.
\textbf{Đáp án của câu hỏi là D}
}
\end{ex}

\begin{ex} %[2H9V47]
Một hộp quà hình lập phương có cạnh $10\text{ cm}$. Hỏi dùng một sợi dây để thắt nơ viền quanh $4$ mặt xung quanh của hộp quà cần ít nhất bao nhiêu xăng-ti-mét dây (không tính phần nơ)?
\choice
{\True $40\text{ cm}$}
{$30\text{ cm}$}
{$50\text{ cm}$}
{$20\text{ cm}$}
\loigiai{
Viền quanh 4 mặt gồm 4 cạnh: $10 \times 4 = 40\text{ (cm)}$.
\textbf{Đáp án của câu hỏi là A}
}
\end{ex}

\begin{ex} %[2H9V47]
Chiếc trống trường học thường có thân trống phình nhẹ nhưng hai mặt trống có dạng phẳng hình tròn. Dạng hình học gần giống nhất của chiếc trống là
\choice
{Khối cầu}
{Khối lập phương}
{\True Khối trụ}
{Khối hộp chữ nhật}
\loigiai{
Trống trường có dạng hình khối trụ.
\textbf{Đáp án của câu hỏi là C}
}
\end{ex}

\begin{ex} %[2H9V47]
Có $18$ lon nước ngọt hình trụ đóng đều vào $3$ lốc. Mỗi lốc có bao nhiêu lon nước ngọt?
\choice
{$5$ lon}
{\True $6$ lon}
{$7$ lon}
{$8$ lon}
\loigiai{
$18 : 3 = 6$ (lon).
\textbf{Đáp án của câu hỏi là B}
}
\end{ex}

\begin{ex} %[2H9V47]
Để lăn một vật trên mặt đất một cách dễ dàng và tốn ít sức nhất, người ta nên dùng vật có dạng khối nào?
\choice
{Khối lập phương}
{Khối hộp chữ nhật}
{Khối tam giác}
{\True Khối trụ hoặc khối cầu}
\loigiai{
Bề mặt cong của khối trụ và khối cầu giúp vật lăn tròn dễ dàng.
\textbf{Đáp án của câu hỏi là D}
}
\end{ex}

\begin{ex} %[2H9V47]
Bé xếp một chiếc xe lửa đồ chơi gồm $1$ khối lập phương làm đầu tàu và $5$ khối hộp chữ nhật làm các toa tàu. Chiếc xe lửa có tất cả bao nhiêu khối hình?
\choice
{\True $6$ khối hình}
{$5$ khối hình}
{$7$ khối hình}
{$8$ khối hình}
\loigiai{
$1 + 5 = 6$ (khối hình).
\textbf{Đáp án của câu hỏi là A}
}
\end{ex}

\textbf{Phần 2. Trắc nghiệm đúng sai}
\setcounterpageref{ex}{0}

\begin{ex} %[2H9V47]
Xét tính đúng sai của các bài toán hình khối thực tế:
\choiceTF
{\True Hộp phấn viết bảng dạng khối hộp chữ nhật có 6 mặt, 8 đỉnh và 12 cạnh.}
{Một khối lập phương có thể lăn tròn về mọi phía như quả bóng.}
{\True Chiếc bánh chưng ngày Tết có hình dạng gần giống khối hộp chữ nhật hoặc khối lập phương dẹt.}
{Tất cả các mặt của khối hộp chữ nhật đều bắt buộc phải là hình vuông.}
\loigiai{
\begin{itemize}
    \item Ý a) Đúng: Đặc điểm cơ bản của hình hộp chữ nhật. Đúng.
    \item Ý b) Sai: Khối lập phương có các mặt phẳng và góc cạnh nên không lăn tròn được.
    \item Ý c) Đúng: Đúng.
    \item Ý d) Sai: Các mặt thường là hình chữ nhật (nếu là hình vuông hết thì thành khối lập phương).
\end{itemize}
\begin{itemize}
    \item a) \textbf{Đúng}.
    \item b) \textbf{Sai}.
    \item c) \textbf{Đúng}.
    \item d) \textbf{Sai}.
\end{itemize}
}
\end{ex}

\begin{ex} %[2H9V47]
Một kệ đồ chơi có: 6 quả bóng tròn, 4 chiếc trống hình trụ, 5 khối rubik lập phương:
\choiceTF
{\True Kệ có tất cả 15 món đồ chơi.}
{\True Số đồ chơi dạng khối cầu nhiều nhất trong ba loại.}
{Số trống hình trụ nhiều hơn số khối rubik lập phương.}
{Nếu lấy bớt đi 2 quả bóng thì số bóng ít hơn số trống.}
\loigiai{
\begin{itemize}
    \item Ý a) Đúng: $6 + 4 + 5 = 15$ món đồ chơi.
    \item Ý b) Đúng: Bóng tròn có 6 quả, nhiều nhất. Đúng.
    \item Ý c) Sai: Trống có 4, rubik có 5 nên trống ít hơn rubik.
    \item Ý d) Sai: Bóng còn: $6 - 2 = 4$, bằng số trống (không phải ít hơn).
\end{itemize}
\begin{itemize}
    \item a) \textbf{Đúng}.
    \item b) \textbf{Đúng}.
    \item c) \textbf{Sai}.
    \item d) \textbf{Sai}.
\end{itemize}
}
\end{ex}

\textbf{Phần 3. Trả lời ngắn}
\setcounterpageref{ex}{0}

\begin{ex} %[2H9V47]
Một khối lập phương có mấy mặt?
\shortans{6}
\loigiai{
\textbf{Đáp số:} $6$\\
Khối lập phương có 6 mặt.
}
\end{ex}

\begin{ex} %[2H9V47]
Có 24 lon sữa hình trụ đóng đều vào 4 hộp carton. Hỏi mỗi hộp có bao nhiêu lon sữa?
\shortans{6}
\loigiai{
\textbf{Đáp số:} $6$\\
$24 : 4 = 6$ lon.
}
\end{ex}

\begin{ex} %[2H9V47]
Bác thợ mộc cưa một khúc gỗ hình trụ dài 20 dm thành 4 khúc bằng nhau. Mỗi khúc dài bao nhiêu đề-xi-mét?
\shortans{5}
\loigiai{
\textbf{Đáp số:} $5$\\
$20 : 4 = 5\text{ dm}$.
}
\end{ex}

\begin{ex} %[2H9V47]
Xếp các khối lập phương nhỏ thành một hàng dài gồm 9 khối, mỗi khối cạnh 3 cm. Hỏi hàng đó dài bao nhiêu xăng-ti-mét?
\shortans{27}
\loigiai{
\textbf{Đáp số:} $27$\\
$3 \times 9 = 27\text{ cm}$.
}
\end{ex}

\begin{ex} %[2H9V47]
Một hộp đồ chơi có 12 khối cầu và 18 khối trụ. Hỏi hộp đó có tất cả bao nhiêu khối hình?
\shortans{30}
\loigiai{
\textbf{Đáp số:} $30$\\
$12 + 18 = 30$ khối.
}
\end{ex}

\begin{ex} %[2H9V47]
Có bao nhiêu đỉnh trên một khối lập phương?
\shortans{8}
\loigiai{
\textbf{Đáp số:} $8$\\
Khối lập phương có 8 đỉnh.
}
\end{ex}
